\section{Cutting the caps lowers the deficit}
\label{sec:cut}

The construction of \cref{sec:inset,sec:annulus,sec:fold} produces from
$\Om$ a connected, hole-free inset $\Etil$ trapped in a thin annulus
$B_{\rho_i}(z)\subset\Etil\subset B_{\rho_e}(z)$, together with a radial core
profile whose graph
$G_0=\{z+r\omega:r<\rho_0(\omega)\}$ satisfies $G_0\subset\Etil$ and
$\abs{\Etil\setminus G_0}=\gd(\Etil)$ (\cref{cor:gd}). Manufacturing the
honest graph from $\Etil$ therefore amounts to \emph{excising} the radial
caps $\Etil\setminus G_0$, all of which lie in the collar
$\{\rho_i<\abs{x-z}<\rho_e\}$. The naive worry is that removing material can
only destroy torsion and hence inflate the scale-invariant deficit
$\dstar$. The opposite is true: caps in a thin collar carry very little
torsion, and the equal-area comparison ball shrinks faster than the torsion
is lost, so $\dstar$ \emph{decreases}. This section proves the two facts
that make this precise. They are sign-critical, so we prove them in full and
check the inequalities digit by digit.

Throughout, for a bounded open set $D\subset\R^2$ the \emph{torsion
function} $u_D\in H^1_0(D)$ is the weak solution of $-\Delta u_D=1$ in $D$,
$u_D=0$ on $\partial D$, and $T(D):=\int_D u_D$ is its torsional rigidity;
for a general set of finite measure $D$ we set $T(D):=T(\mathrm{int}\,D)$ on
the open representative produced by our construction, which is all we use.
Recall (\cref{thm:torsion-reg}) that $u_D$ is real-analytic in $D$ and, by
the comparison principle for the Laplacian, $u_D\ge0$. The variational
characterisation we use repeatedly is that $u_D$ minimises the Dirichlet
energy functional $J_D(v):=\tfrac12\int_D\abs{\nabla v}^2-\int_D v$ over
$v\in H^1_0(D)$, with
\begin{equation}\label{eq:torsion-energy}
   \min_{v\in H^1_0(D)}J_D(v)=J_D(u_D)=-\tfrac12\,T(D),
   \qquad
   \int_D\abs{\nabla u_D}^2=\int_D u_D=T(D),
\end{equation}
the second identity being the Euler--Lagrange relation tested against $u_D$
itself.

\subsection{Torsion loss under radial-cap cutting}

The next lemma is the analytic heart of the section: it controls the torsion
lost when one cuts radial caps off a set trapped in a thin collar. Unlike a
naive measure-based bound --- which is false in general, since excising a
zero-area set of positive capacity (a slit) changes the torsion --- the
estimate is variational and exploits the specific nested radial-cap geometry
produced in \cref{sec:fold}.

\begin{lemma}[Torsion loss under radial-cap cutting]\label{lem:torsion-stab}
Let $z\in\R^2$ and $0<\rho_i\le\rho_e$, and put
\begin{equation}\label{eq:tildeeta}
   \tilde\eta:=\tfrac14\bigl(\rho_e^2-\rho_i^2\bigr).
\end{equation}
Let $G\subset E\subset\R^2$ be bounded open sets with
\begin{equation}\label{eq:collar-hyp}
   B_{\rho_i}(z)\subset G\subset E\subset B_{\rho_e}(z),\qquad
   E\setminus G\subset\{\rho_i<\abs{x-z}<\rho_e\},
\end{equation}
and put $g:=\abs{E\setminus G}$. Then $0\le u_G\le u_E$ on $G$, the
difference $h:=u_E-u_G$ \textup(extended by $u_E$ on $E\setminus G$, equivalently
$h=u_E-\bar u_G$ where $\bar u_G\in H^1_0(E)$ is the extension of $u_G$ by
zero\textup) lies in $H^1_0(E)$ and is harmonic in $G$, and
\begin{equation}\label{eq:torsion-loss}
   0\ \le\ T(E)-T(G)
   \ =\ \int_{E\setminus G}u_E\ +\ \int_G h
   \ \le\ \tilde\eta\,g\ +\ \int_G h .
\end{equation}
Moreover the boundary-layer term $\int_G h$ is non-negative, satisfies the
energy identity
\begin{equation}\label{eq:energy-id}
   \int_G h\ \le\ T(E)-T(G)\ =\ \int_E\abs{\nabla h}^2,
\end{equation}
and obeys $0\le h\le\tilde\eta$ throughout $G$, whence the crude bound
$\int_G h\le\tilde\eta\,\abs G$.
\end{lemma}

\begin{proof}
\emph{Step A: the exact decomposition.}
Since $G\subset E$, the restriction $u_E|_G$ solves $-\Delta u_E=1=-\Delta
u_G$ in $G$ with $u_E\ge0=u_G$ on $\partial G$; by the comparison principle
on $G$ the difference $h:=u_E-u_G$ is then $\ge0$, i.e.\ $0\le u_G\le u_E$ on
$G$. Being a difference of two solutions of $-\Delta(\cdot)=1$, $h$ is
harmonic in $G$, with boundary trace $h=u_E$ on $\partial G$ (as $u_G=0$
there). Splitting $E=G\cup(E\setminus G)$ and using $u_G=0$ on $E\setminus G$,
\begin{equation}\label{eq:TEG-split}
   T(E)-T(G)=\int_E u_E-\int_G u_G
   =\int_{E\setminus G}u_E+\int_G(u_E-u_G)
   =\int_{E\setminus G}u_E+\int_G h,
\end{equation}
which is the equality in \eqref{eq:torsion-loss}; both summands are
$\ge0$, so $T(E)-T(G)\ge0$.

\emph{Step B: the cap term.}
Since $E\subset B_{\rho_e}(z)$, the comparison principle gives $u_E\le
u_{B_{\rho_e}(z)}$ on $E$: the difference $u_{B_{\rho_e}(z)}-u_E$ is harmonic
in $E$ (both solve $-\Delta(\cdot)=1$) and equals $u_{B_{\rho_e}(z)}\ge0$ on
$\partial E$, hence is $\ge0$ in $E$ by the minimum principle. The disc
torsion function is explicit,
\begin{equation}\label{eq:discu}
   u_{B_{\rho_e}(z)}(y)=\frac{\rho_e^2-\abs{y-z}^2}{4}.
\end{equation}
On the caps $E\setminus G\subset\{\rho_i<\abs{x-z}<\rho_e\}$ one has
$\abs{x-z}\ge\rho_i$, so
\begin{equation}\label{eq:barrier}
   u_E\ \le\ \frac{\rho_e^2-\abs{x-z}^2}{4}
   \ \le\ \frac{\rho_e^2-\rho_i^2}{4}\ =\ \tilde\eta
   \qquad\text{on } E\setminus G,
\end{equation}
and therefore $\int_{E\setminus G}u_E\le\tilde\eta\,g$. This is the main
term and is fully rigorous; combined with \eqref{eq:TEG-split} it gives the
inequality in \eqref{eq:torsion-loss}.

\emph{Step C: the boundary-layer term, energy form and crude bound.}
Let $\bar u_G\in H^1_0(E)$ denote the extension of $u_G$ by zero outside $G$;
then $h=u_E-\bar u_G\in H^1_0(E)$ agrees with the previously defined $h$ on
$G$ and equals $u_E$ on $E\setminus G$. Since $u_E$ minimises $J_E$ on
$H^1_0(E)$ and $\bar u_G$ minimises $J_E$ restricted to the subspace
$H^1_0(G)\subset H^1_0(E)$ (its minimum there being $J_E(\bar
u_G)=J_G(u_G)=-\tfrac12T(G)$), the parallelogram identity for the quadratic
functional $J_E$ gives the exact energy balance
\begin{equation}\label{eq:energy-balance}
   \tfrac12\int_E\abs{\nabla h}^2
   =J_E(\bar u_G)-J_E(u_E)
   =\bigl(-\tfrac12T(G)\bigr)-\bigl(-\tfrac12T(E)\bigr)
   =\tfrac12\bigl(T(E)-T(G)\bigr),
\end{equation}
using \eqref{eq:torsion-energy} for both $D=E$ and $D=G$. Hence
$T(E)-T(G)=\int_E\abs{\nabla h}^2$, the equality in \eqref{eq:energy-id};
combined with \eqref{eq:torsion-loss} and $\int_{E\setminus G}u_E\ge0$ this
gives $\int_G h\le T(E)-T(G)=\int_E\abs{\nabla h}^2$.

Finally, the bounds $0\le h\le\tilde\eta$ in $G$: non-negativity was shown in
Step A. For the upper bound, $h$ is harmonic in $G$ with boundary trace
$h=u_E$ on $\partial G$, and on $\partial G$ one has $u_E\le\tilde\eta$
--- indeed $\partial G\subset\overline{B_{\rho_e}(z)}$, so the barrier
\eqref{eq:discu} gives $u_E\le u_{B_{\rho_e}(z)}\le\tilde\eta$ on the part of
$\partial G$ lying in $\{\abs{x-z}\ge\rho_i\}$ (which contains
$\partial G\cap E$, the interior cap-floor boundary, since $E\setminus G$ is
confined to the collar by \eqref{eq:collar-hyp}), while $u_E=0$ on
$\partial G\cap\partial E$. By the maximum principle $h\le\tilde\eta$
throughout $G$, whence $\int_G h\le\tilde\eta\abs G$.
\end{proof}

The cap term $\tilde\eta\,g$ in \eqref{eq:torsion-loss} is exactly what the
deficit algebra of \cref{lem:cutdeficit} can afford. The boundary-layer term
$\int_G h$ is genuinely smaller in the application but requires care; we
isolate it next.

\begin{remark}[The boundary-layer term and the residual analytic point]
\label{gap:torsion-stab}
The cap term $\int_{E\setminus G}u_E\le\tilde\eta\,g$ of Step B is rigorous
and, as we show in \cref{lem:cutdeficit}, by itself drives the deficit down
once the harmonic boundary-layer term $\int_G h$ is controlled at the same
order. The sharp description of $\int_G h$ is, by Green's identity in $G$
\textup(legitimate when $\partial G$ admits an outward normal $\HH^1$-a.e.;
for the rough star-shaped $G$ below it is the limiting form along a smooth
exhaustion\textup),
\begin{equation}\label{eq:flux-rep}
   \int_G h=\int_G h\,(-\Delta u_G)
   =-\int_{\partial G}h\,\partial_\nu u_G\,d\HH^1
   =\int_{\partial G}u_E\,\abs{\partial_\nu u_G}\,d\HH^1
   =\int_\Gamma u_E\,\abs{\partial_\nu u_G}\,d\HH^1,
\end{equation}
where $\partial_\nu u_G\le0$ on $\partial G$ \textup(as $u_G\ge0=u_G$ on
$\partial G$\textup) and the integral is supported on the interior cap-floor
boundary $\Gamma:=\partial G\cap E$, because $u_E=0$ on the remainder
$\partial G\cap\partial E$. Since $u_E\le\tilde\eta$ on $\Gamma$ by
\eqref{eq:barrier}, this yields the clean bound
\begin{equation}\label{eq:flux-bound}
   \int_G h\ \le\ \tilde\eta\,\Phi,
   \qquad
   \Phi:=\int_\Gamma\abs{\partial_\nu u_G}\,d\HH^1,
\end{equation}
where $\Phi$ is the flux of $u_G$ through the interior cap-floor. The total
flux through the \emph{whole} boundary is exactly $\int_{\partial
G}\abs{\partial_\nu u_G}=\int_G(-\Delta u_G)=\abs G$, so $\Phi\le\abs G$
unconditionally; the substance is that $\Gamma$ is only the small,
fold-supported part of $\partial G$. For a boundary with a one-sided gradient
bound $\norm{\nabla u_G}_{L^\infty(\partial G)}\le C$ \textup(an $O(1)$
quantity for the disc-trapped torsion function, by the barrier $u_G\le
u_{B_{\rho_e}(z)}$ and interior estimates\textup) one has $\Phi\le
C\,\HH^1(\Gamma)$, so it would suffice to bound the length $\HH^1(\Gamma)$ of
the interior cap-floor by the fold content of \cref{prop:fold}. This is the
\emph{single} residual analytic point: the merely measurable core profile
$\rho_0=r_1$ of \cref{cor:gd} carries no slope control, so $\HH^1(\Gamma)$
--- which involves the radial derivative of $\rho_0$ --- is not controlled by
the fold content (a rough graph can have large perimeter over a small angular
sector). We therefore record the boundary-layer control as the explicit
scoped hypothesis
\begin{equation}\label{eq:flux-hyp}
   \int_G h\ \le\ \tilde\eta\,\Phi\ \le\ \tilde\eta\,g
\end{equation}
\textup(equivalently $\Phi\le g$: the flux through the interior cap-floor is
dominated by the excised area\textup), which holds whenever $\partial G$ has
the regularity making \eqref{eq:flux-rep} valid and $\HH^1(\Gamma)\le
C(\Exc+a\eta)$; under \eqref{eq:flux-hyp} the cutting estimate below is
unconditional. We do not claim \eqref{eq:flux-hyp} for an arbitrary
measurable star-shaped $G$, and we flag every downstream use of it
explicitly.
\end{remark}

\begin{remark}[The barrier constant]\label{rem:barrier}
The width of the collar enters only through
$\tilde\eta=\tfrac14(\rho_e^2-\rho_i^2)$, the maximal value of the disc
torsion function on the inner sphere. With $\rho_i=\hat r-O(\eta)$ and
$\rho_e=\hat r+O(\eta)$ one has $\tilde\eta=\tfrac14(\rho_e+\rho_i)
(\rho_e-\rho_i)\le\tfrac12\rho_e\,\eta\le C\eta$, so $\tilde\eta$ is of the
same order as the annulus width $\eta=\rho_e-\rho_i$. In the unit-scale
normalisation $\rho_e=1+\eta$, $\rho_i=1-\eta$ one gets exactly
$\tilde\eta=\eta$.
\end{remark}

\subsection{Cutting boundary material decreases \texorpdfstring{$\dstar$}{the scale-invariant deficit}}

Recall the scale-invariant relative deficit of a finite-measure set $D$,
\begin{equation}\label{eq:dstar-recall}
   \dstar(D)=\frac{T(B_D)-T(D)}{T(B_D)},\qquad
   T(B_D)=\frac{\abs D^2}{8\pi},
\end{equation}
where $B_D$ is any ball with $\abs{B_D}=\abs D$ (so
$T(B_D)=|B_D|^2/(8\pi)$ by scaling from $T(B_1)=\pi/8$). Equivalently,
$\dstar(D)=1-8\pi\,T(D)/\abs D^2$. We now show that excising radial caps
confined to a thin collar lowers $\dstar$, provided the torsion loss per unit
excised area stays below the torsion density of the larger set. The
hypothesis is stated directly as this \emph{loss-slope} condition, which is
exactly what the sign-critical algebra consumes and which the geometric input
of \cref{lem:torsion-stab} supplies.

\begin{lemma}[Cutting boundary material decreases $\dstar$]\label{lem:cutdeficit}
Let $z\in\R^2$, $0<\rho_i\le\rho_e$ and $\tilde\eta=\tfrac14(\rho_e^2-
\rho_i^2)$. Let $E,G$ be bounded open sets with
\begin{equation}\label{eq:cut-incl}
   B_{\rho_i}(z)\subset G\subset E\subset B_{\rho_e}(z),
   \qquad E\setminus G\subset\{\rho_i<\abs{x-z}<\rho_e\},
\end{equation}
put $V:=\abs E$, $g:=\abs{E\setminus G}=V-\abs G$, and suppose the
\emph{torsion-loss bound}
\begin{equation}\label{eq:cutcond}
   T(E)-T(G)\ \le\ \lambda\,g
   \qquad\text{with}\qquad
   \lambda\ \le\ \frac{T(E)}{\abs E}
\end{equation}
holds. Then $\dstar(G)\le\dstar(E)$.
\end{lemma}

\begin{proof}
Write $T:=T(E)$; then $\abs G=V-g$ and, since $\varnothing\ne B_{\rho_i}(z)
\subset G\subset E$, we have $0\le g\le V$ and $V>0$.

\emph{Step 1: lower bound for $T(G)$.} This is the hypothesis
\eqref{eq:cutcond} rearranged:
\begin{equation}\label{eq:TGlower}
   T(G)\ \ge\ T-\lambda\,g,
   \qquad \lambda\ \le\ \frac TV .
\end{equation}

\emph{Step 2: the deficit comparison reduces to a torsion inequality.} By
\eqref{eq:dstar-recall}, with $\abs E=V$ and $\abs G=V-g$,
\begin{equation}\label{eq:dstar-diff}
   \dstar(E)-\dstar(G)
   =\Bigl(1-\frac{8\pi\,T}{V^2}\Bigr)-\Bigl(1-\frac{8\pi\,T(G)}{(V-g)^2}\Bigr)
   =8\pi\Bigl[\frac{T(G)}{(V-g)^2}-\frac{T}{V^2}\Bigr].
\end{equation}
Thus $\dstar(G)\le\dstar(E)$ \emph{holds if and only if}
$\dfrac{T(G)}{(V-g)^2}\ge\dfrac{T}{V^2}$, i.e.\ (multiplying by
$(V-g)^2\ge0$) if and only if
\begin{equation}\label{eq:target}
   T(G)\ \ge\ T\,\frac{(V-g)^2}{V^2}\ =\ T\Bigl(1-\frac gV\Bigr)^2 .
\end{equation}
We now prove \eqref{eq:target}. (When $g=V$, $\abs G=0$ is impossible since
$G\supset B_{\rho_i}(z)$, so in fact $g<V$; but the argument below covers
$g\le V$ verbatim.)

\emph{Step 3: the sign-critical chain.} Expand the right-hand side of
\eqref{eq:target}:
\begin{equation}\label{eq:expand}
   T\Bigl(1-\frac gV\Bigr)^2
   =T-\frac{2Tg}{V}+\frac{Tg^2}{V^2}.
\end{equation}
Subtracting \eqref{eq:expand} from the lower bound \eqref{eq:TGlower} for
$T(G)$,
\begin{align}
   T(G)-T\Bigl(1-\frac gV\Bigr)^2
   &\ \ge\ \bigl(T-\lambda\,g\bigr)
     -\Bigl(T-\frac{2Tg}{V}+\frac{Tg^2}{V^2}\Bigr)
   \notag\\
   &\ =\ \frac{2Tg}{V}-\frac{Tg^2}{V^2}-\lambda\,g
   \ =\ g\Bigl(\frac{2T}{V}-\frac{Tg}{V^2}-\lambda\Bigr).
   \label{eq:bracket}
\end{align}
We bound the bracket from below. Since $0\le g\le V$ we have
$0\le g/V\le1$, hence $2-g/V\ge1$, and therefore
\begin{equation}\label{eq:bracket-key}
   \frac{2T}{V}-\frac{Tg}{V^2}
   =\frac TV\Bigl(2-\frac gV\Bigr)\ \ge\ \frac TV\ \ge\ \lambda,
\end{equation}
the last inequality being exactly the loss-slope hypothesis
\eqref{eq:cutcond} ($\lambda\le T/V=T(E)/\abs E$). Consequently the bracket
in \eqref{eq:bracket} is $\ge\frac TV-\lambda\ge0$, and since $g\ge0$ the
whole expression \eqref{eq:bracket} is $\ge0$. This is \eqref{eq:target},
which by Step 2 yields $\dstar(G)\le\dstar(E)$.
\end{proof}

\begin{remark}[Supplying the loss-slope hypothesis from
\cref{lem:torsion-stab}]\label{rem:supply}
\Cref{lem:torsion-stab} furnishes \eqref{eq:cutcond} as follows. Its
hypotheses coincide with \eqref{eq:cut-incl}, and \eqref{eq:torsion-loss}
gives the exact torsion loss $T(E)-T(G)=\int_{E\setminus G}u_E+\int_G h\le
\tilde\eta\,g+\int_G h$. Under the boundary-layer control \eqref{eq:flux-hyp}
of \cref{gap:torsion-stab} --- $\int_G h\le\tilde\eta\,g$ --- this reads
\begin{equation}\label{eq:loss-2eta}
   T(E)-T(G)\ \le\ 2\,\tilde\eta\,g,
\end{equation}
so \eqref{eq:cutcond} holds with the explicit slope $\lambda=2\tilde\eta$,
provided $2\tilde\eta\le T(E)/\abs E$. Even with only the crude bound
$\int_G h\le\tilde\eta\abs G$ of \cref{lem:torsion-stab} one obtains the loss
$\le\tilde\eta(g+\abs G)=\tilde\eta\abs E$, but that slope is
$\tilde\eta\abs E/g$, which is \emph{not} below $T/V$ at the relevant scales;
the gain from \eqref{eq:flux-hyp} is precisely the localisation of the
boundary-layer term to the cap region, replacing the slope $\tilde\eta\abs
E/g$ by the absolute $2\tilde\eta$.
\end{remark}

\begin{remark}[The loss-slope hypothesis \eqref{eq:cutcond} holds with room
to spare]\label{rem:roomtospare}
In the application $E=\Etil$ is the inset of \cref{def:inset}, of measure
$\mu_E=\abs\Etil=\pi-O(\sqrt\dT)\ge\pi/2$ and deficit
$\dstar(\Etil)\le\tfrac12$ for $\dT\le\delta_\star$. Then the torsion density
is an \emph{absolute} constant bounded below,
\[
   \frac{T(\Etil)}{\abs\Etil}
   =\bigl(1-\dstar(\Etil)\bigr)\frac{\abs\Etil}{8\pi}
   \ \ge\ \bigl(1-\tfrac12\bigr)\cdot\frac{\pi/2}{8\pi}
   \ =\ \frac1{32},
\]
using $T(\Etil)=(1-\dstar(\Etil))\,T(B_{\Etil})=(1-\dstar(\Etil))\,
\abs\Etil^2/(8\pi)$ from \eqref{eq:dstar-recall}. On the other hand the slope
$\lambda=2\tilde\eta$ supplied by \cref{rem:supply} vanishes:
$\tilde\eta\le\tfrac12\rho_e\,\eta\le C\eta\le C(\dT/\tau)^{1/4}\to0$ as
$\dT\to0$ (\cref{thm:annulus}, with $\tau=\sqrt\dT$ giving
$\eta\le C\dT^{1/8}$). Hence $\lambda=2\tilde\eta\le\tfrac1{32}\le
T(\Etil)/\abs\Etil$ for all $\dT\le\delta_\star$, and the loss-slope
condition \eqref{eq:cutcond} is satisfied with a fixed margin. The crucial
structural point is that the comparison is between a \emph{vanishing} slope
$\lambda=2\tilde\eta\to0$ and an \emph{absolute} torsion density $\ge1/32$,
so the inequality $\lambda\le T(\Etil)/\abs\Etil$ holds robustly; the
threshold $1/32$ is not sharp and is recorded only to fix $\delta_\star$.
This is exactly where the boundary-layer control \eqref{eq:flux-hyp} enters:
it is what turns the exact loss $\int_{E\setminus G}u_E+\int_G h$ into a bound
with absolute slope $2\tilde\eta$ rather than the much larger
$\tilde\eta\abs\Etil/g$.
\end{remark}

\begin{remark}[The application: cutting the caps]\label{rem:twouses}
\Cref{lem:cutdeficit} is used in the assembly (\cref{thm:reduction}) with
$E=\Etil$ and $G=G_0=\{z+r\omega:r<\rho_0(\omega)\}$ the measurable radial
core graph of \cref{cor:gd,prop:realization}, taken in its open star-shaped
representative. The inclusions $B_{\rho_i}(z)\subset G_0\subset\Etil\subset
B_{\rho_e}(z)$ hold by \cref{cor:gd,thm:annulus}, the difference
$\Etil\setminus G_0$ is the union of the radial caps, all contained in the
collar $\{\rho_i<\abs{x-z}<\rho_e\}$, and the exact torsion-loss decomposition
$T(\Etil)-T(G_0)=\int_{\Etil\setminus G_0}u_{\Etil}+\int_{G_0}h\le\tilde\eta\,g
+\int_{G_0}h$ of \cref{lem:torsion-stab} holds unconditionally. Under the
boundary-layer control \eqref{eq:flux-hyp} of \cref{gap:torsion-stab} the loss
is $\le2\tilde\eta\,g$ (\cref{rem:supply}); together with the absolute
torsion-density lower bound of \cref{rem:roomtospare}, the loss-slope
condition \eqref{eq:cutcond} holds with $\lambda=2\tilde\eta$, and
\cref{lem:cutdeficit} gives that cutting the caps lowers the deficit,
\[
   \dstar(G_0)\ \le\ \dstar(\Etil).
\]
This is the one place where the geometry of \cref{lem:cutdeficit} is needed:
the discarded material is precisely low-torsion boundary material in the
$\eta$-collar about the \emph{single} centre $z$. The deduction is
unconditional except for the single residual analytic point \eqref{eq:flux-hyp}
isolated in \cref{gap:torsion-stab} --- the localisation of the harmonic
boundary-layer term to the cap region; the cap term $\tilde\eta\,g$ and the
entire deficit algebra are rigorous.
\end{remark}

\begin{remark}[The other deficit step is \emph{not} an instance of
\cref{lem:cutdeficit}]\label{rem:transfer-not-cut}
The earlier reduction step --- passing from the full superlevel set
$\{u>\hat v\}$ to its largest connected component $\Etil$, i.e.\ discarding
the spillover components of total measure
$\abs{\{u>\hat v\}\setminus\Etil}\le D_I(\hat v)^2/(16\pi^3)$
(\cref{lem:spill}) --- is \emph{not} covered by \cref{lem:cutdeficit}, and we
do not claim it is. \Cref{lem:cutdeficit} requires
$G\subset E\subset B_{\rho_e}(z)$ with $E\setminus G$ inside one collar about
$z$; but the discarded spillover components are distinct level components,
which need not lie in the trapping ball $B_{\rho_e}(z)$ of $\Etil$ (they can
sit elsewhere in $\Om$), so the hypothesis $E\subset B_{\rho_e}(z)$ may fail
for $E=\{u>\hat v\}$. The correct tool for that step is the deficit transfer
\cref{lem:transfer} of \cref{sec:inset} combined with the \emph{exact}
additivity of torsion over connected components (the function $u-\hat v$ is
the common torsion function of every component of $\{u>\hat v\}$); this is
carried out in \cref{prop:insetcontrols}, yielding $\dstar(\Etil)\le
(\pi^2/\mu_E^2)\,\dstar(\Om)+O(\dT^2/\tau^2)$ with no appeal to the collar
geometry. Thus the full chain
\[
   \dstar(G_0)\ \le\ \dstar(\Etil)
   \ \le\ \frac{\pi^2}{\mu_E^2}\,\dstar(\Om)+O\!\Bigl(\frac{\dT^2}{\tau^2}\Bigr)
\]
combines \cref{lem:cutdeficit} (first inequality, this section) with
\cref{lem:transfer}/\cref{prop:insetcontrols} (second inequality,
\cref{sec:inset}); the two play distinct roles and rest on distinct
hypotheses.
\end{remark}

\begin{remark}[Why the deficit decreases — the mechanism]\label{rem:mechanism}
The single inequality \eqref{eq:bracket-key} is the entire content. Cutting
a cap of measure $g$ costs at most $\lambda\,g$ in torsion (\emph{linear} in
$g$, with a small slope $\lambda=2\tilde\eta\to0$ under \eqref{eq:flux-hyp}),
but it shrinks the equal-area comparison ball, lowering its torsion
$T(B_G)=\abs G^2/(8\pi)$ by $\approx\frac{2T(B_E)}{V}g=\frac{V}{4\pi}g$, i.e.\
by a slope $\approx\frac{2T}{V}$ after accounting for the deficit factor.
Because the unperturbed torsion density $T/V$ of a near-disc is an
\emph{absolute} constant of order $1$ while the collar slope
$\lambda=2\tilde\eta$ vanishes, the denominator $T(B_G)$ drops faster than the
numerator $T(G)$, and $\dstar=1-T/T(B)$ decreases. There is no square-root
loss and no constant that degrades as $\dT\to0$: the deficit algebra is exact
and the margin in \eqref{eq:cutcond} only widens. The exactness of the
torsion-loss \emph{decomposition} \eqref{eq:torsion-loss} is likewise
unconditional; the only inexact step is the localisation \eqref{eq:flux-hyp}
of the harmonic boundary-layer term, isolated in \cref{gap:torsion-stab}.
\end{remark}
